\par\indent\SourceLeaf{236}%
III. Statuendo \(x=r\cos\varphi,\ y=r\sin\varphi\), fit
\[
z=\tfrac{1}{2}(T\cos\varphi^2+V\sin\varphi^2)rr+\Omega
\]
unde colligitur, si sectio fiat per planum superficiei in \(A\) normale et cum axe ipsarum \(x\) angulum \(\varphi\) efficiens, oriri curvam planam, cuius radius curvaturae in puncto \(A\) sit
\[
=\frac{1}{T\cos\varphi^2+V\sin\varphi^2}
\]
\par\indent IV. Quoties itaque habetur \(T=V\), radii curvaturae in \textit{cunctis} planis normalibus aequales erunt. Si vero \(T\) et \(V\) sunt inaequales, manifestum est, quum \(T\cos\varphi^2+V\sin\varphi^2\) pro quovis valore anguli \(\varphi\) cadat intra \(T\) et \(V\), radios curvaturae in sectionibus principalibus, in I et II consideratis, referri ad curvaturas extremas, puta alterum ad curvaturam maximam, alterum ad minimam, si \(T\) et \(V\) eodem signo affectae sint, contra alterum ad maximam convexitatem, alterum ad maximam concavitatem, si \(T\) et \(V\) signis oppositis gaudeant. Hae conclusiones omnia fere continent, quae ill. \textsc{Euler} de curvatura superficierum curvarum primus docuit.

V. Mensura curvaturae superficiei curvae in puncto \(A\) autem nanciscitur expressionem simplicissimam \(k=TV\), unde habemus

\textsc{Theorema.} \textit{Mensura curvaturae in quovis superficiei puncto aequalis est fractioni, cuius numerator unitas, denominator autem productum duorum radiorum curvaturae extremorum in sectionibus per plana normalia.}

Simul patet, mensuram curvaturae fieri positivam pro superficiebus concavo-concavis vel convexo-convexis (quod discrimen non est essentiale), negativam vero pro concavo-convexis. Si superficies constat e partibus utriusque generis, in earum confiniis mensura curvaturae evanescens esse debebit. De indole superficierum curvarum talium, in quibus mensura curvaturae ubique evanescit, infra pluribus agetur.

\Article{9}
Formula generalis pro mensura curvaturae in fine art. 7 proposita, omnium simplicissima est, quippe quae quinque tantum elementa implicat; ad magis complicatam, scilicet novem elementa involventem, deferimur, si adhibere volumus
